\documentclass[a4paper,11pt]{article}
\usepackage[utf8]{inputenc}
\usepackage[T1]{fontenc}
\usepackage[french]{babel}
\usepackage[left=2cm,right=2cm,top=2cm,bottom=2cm]{geometry}
\usepackage{xcolor}
\usepackage{hyperref}
\usepackage{fontawesome5}
\usepackage{parskip}

% --- Couleurs EDF ---
\definecolor{primary}{RGB}{0, 90, 156}
\hypersetup{colorlinks=true, urlcolor=primary}

\begin{document}

% --- EXPÉDITEUR ---
\noindent
\begin{minipage}[t]{0.5\textwidth}
    \textbf{Loïc Aron MBASSI EWOLO}\\
    Élève-Ingénieur Bac+5 -- Double diplôme ENSEA\\[3pt]
    \faMapMarker\ Cergy, France\\
    \faPhone\ (+33) 7 59 52 33 98\\
    \faEnvelope\ \href{mailto:loicaron.mbassiewolo@ensea.fr}{loicaron.mbassiewolo@ensea.fr}
\end{minipage}
\hfill
\begin{minipage}[t]{0.4\textwidth}
    \raggedleft
    Cergy, le 5 mars 2026
\end{minipage}

\vspace{1.2cm}

\hfill
\begin{minipage}[t]{0.5\textwidth}
    \raggedleft
    \textbf{EDF}\\
    Service Recrutement\\
    Grenoble, France
\end{minipage}

\vspace{1cm}

\noindent\textbf{Objet :} Candidature -- Alternance Développeur Logiciel Embarqué (F/H)

\vspace{0.8cm}

Madame, Monsieur,

Concevoir un algorithme de SLAM sur Nvidia Jetson pour qu'un robot militaire cartographie un environnement inconnu en temps réel, intégrer des données Lidar 3D et caméra de profondeur, écrire du C/C++ en tenant compte des contraintes mémoire et temps réel : voilà le quotidien du challenge CoHoMa, projet de robotique autonome mené à l'ENSEA en collaboration avec l'armée française. C'est ce contexte de développement embarqué exigeant qui m'a conduit à postuler à l'alternance Développeur Logiciel Embarqué chez EDF.

Élève-ingénieur en double diplôme ENSEA -- École Polytechnique de Yaoundé (Mathématiques Appliquées, Génie Logiciel), je suis actuellement en stage de recherche à l'INRIA Saclay (Équipe TRIBE, campus Institut Polytechnique de Paris), où je développe en Python des pipelines d'analyse statique de code à grande échelle, en pratiquant notamment l'inspection automatisée des comportements réels des SDKs logiciels. Ce travail d'analyse statique et de qualification logicielle est directement en lien avec les exigences de développement sûr que vous pratiquez chez EDF pour les logiciels de sûreté nucléaire.

Sur le plan du développement embarqué à proprement parler, le projet CoHoMa m'a amené à travailler sur des problématiques de fiabilité et de robustesse proches des vôtres : détection de pannes, comportements dégradés, contraintes temps réel strictes. J'ai également mené un projet académique de dix semaines sur la Fault Detection and Fault-Tolerant Control d'un drone (MATLAB/Simulink), portant sur la détection et l'isolation de défauts capteurs et actionneurs, et la reconfiguration de contrôleur en situation dégradée. En parallèle, le projet Harmony Gloves (prototypage de gants instrumentés, STM32, capteurs IMU/Flex, soudure) et le projet TAXI (simulation multiprocessus en C pur, IPC/mémoire partagée, signaux POSIX, Linux) complètent mon expérience en développement bas niveau et en gestion de la concurrence.

EDF est l'un des rares environnements industriels où les exigences logicielles atteignent le niveau de criticité rencontré dans l'embarqué militaire ou spatial, et c'est précisément ce niveau d'exigence qui m'attire. Je suis disponible dès septembre 2026 pour une alternance et je serais heureux d'en discuter lors d'un entretien.

Je reste à votre disposition pour tout complément d'information.

Je vous prie d'agréer, Madame, Monsieur, l'expression de mes salutations distinguées.

\vspace{0.8cm}

\textbf{Loïc Aron MBASSI EWOLO}\\
\textit{Élève-Ingénieur ENSEA / Polytechnique Yaoundé}

\end{document}
